%----------------Cap_02----------------%
\chapter{Fundamentação Teórica e Estado da Arte}
\label{cap:fundamentacao}

Este capítulo estabelece a base teórica e conceitual que fundamenta o desenvolvimento deste trabalho, cobrindo os princípios do controle clássico, a formulação do controle preditivo não-linear, os métodos de identificação de sistemas dinâmicos e as técnicas de inteligência computacional embarcada (\textit{TinyML}) empregadas para a aproximação de políticas ótimas. Ao final, revisam-se os trabalhos correlatos recentes na literatura e evidencia-se a lacuna científica e tecnológica abordada pela pesquisa.

% ============================================================
\section{Controle Clássico}
\label{sec:fund_classico}
% ============================================================

Controladores clássicos como o PID (Proporcional-Integral-Derivativo) constituem a espinha dorsal de inúmeras aplicações industriais e robóticas. O desenvolvimento formal dessa lei de controle, também conhecida como controle de três termos, foi apresentado em 1922 pelo engenheiro Nicolas Minorsky \cite{bennett1996}.

Em plataformas ágeis, como os Veículos Aéreos Não Tripulados (VANTs), o controle PID muitas vezes não é suficiente de forma isolada, sendo frequentemente relegado a uma malha interna rápida de estabilização de atitude ou de velocidade angular. Essa limitação decorre de sua incapacidade intrínseca de considerar restrições operacionais explícitas nos atuadores (como saturações de empuxo) e de antecipar dinâmicas futuras não-lineares, criando a necessidade de estratégias preditivas mais avançadas capazes de gerenciar sistemas com dinâmicas acopladas e restrições físicas rígidas \cite{nguyen2020}.

% ============================================================
\section{Controle Preditivo Não-Linear (NMPC)}
\label{sec:fund_nmpc}
% ============================================================

O Controle Preditivo Não-Linear (NMPC) fundamenta-se na resolução repetida, a cada período de amostragem, de um problema de controle ótimo em tempo discreto e horizonte finito. Um otimizador numérico propaga um modelo preditivo não-linear da planta para calcular uma sequência de controle que minimiza uma função de custo quadrática acumulada, respeitando ativamente restrições de estado e de entrada \cite{mayne2000, rawlings2017}.

Uma das maiores vantagens do NMPC em relação às malhas lineares clássicas reside na característica de \textit{preview} (visão antecipada da referência futura): ao conhecer a trajetória desejada ao longo do horizonte $N$, o controlador inicia a aceleração ou frenagem de maneira preventiva antes que o erro de rastreamento se acumule. Todavia, a resolução iterativa online de um Problema de Programação Não-Linear (NLP) impõe elevado custo computacional, tornando sua aplicação direta em microcontroladores de baixo custo um desafio expressivo em sistemas com dinâmicas rápidas \cite{nguyen2020}.

% ============================================================
\section{MPC Explícito e suas Limitações}
\label{sec:fund_empc}
% ============================================================

Buscando contornar a sobrecarga computacional da otimização numérica online, propôs-se o conceito de MPC Explícito (EMPC). Nessa formulação, o problema de controle ótimo sob restrições lineares é resolvido parametricamente offline por meio de programação multiparamétrica. A lei de controle resultante é expressa como uma função afim por partes (\textit{Piecewise Affine} - PWA) definida sobre partições politópicas do espaço de estados, de modo que a avaliação online restringe-se à busca pontual do polítopo ativo e à avaliação de uma equação linear simples \cite{bemporad2002}.

Apesar de eliminar os resolvedores iterativos no hardware embarcado, o EMPC sofre de um severo problema de escalabilidade dimensional: o número de partições politópicas cresce exponencialmente com o horizonte de predição e com a quantidade de restrições ativas. Essa explosão combinatória inviabiliza o armazenamento da tabela de polítopos na memória estática de microcontroladores para sistemas de ordem média ou com horizontes estendidos \cite{bemporad2002}. Tal limitação estimula a exploração de técnicas de aprendizado de máquina como aproximadores de função com consumo de memória determinístico \cite{alsmeier2024}.

% ============================================================
\section{Aproximação de Leis de Controle por Redes Neurais}
\label{sec:fund_ann_mpc}
% ============================================================

A necessidade de contornar as restrições de memória do EMPC impulsiona o emprego de Redes Neurais Artificiais para aprender a política ótima de controle. Formalmente, \citeonline{karg2020} demonstraram que Redes Neurais Profundas (DNNs) com funções de ativação do tipo ReLU (\textit{Rectified Linear Unit}) são estruturalmente capazes de representar de forma exata qualquer função afim por partes gerada por controladores preditivos lineares e aproximar funções contínuas arbitrárias decorrentes de NMPCs com alta fidelidade.

O desafio central dessa abordagem consiste no controle do erro de aproximação e na manutenção de salvaguardas de estabilidade. Para mitigar esses riscos, investigam-se abordagens estatísticas baseadas na Desigualdade de Hoeffding para certificar a satisfação de restrições com alta probabilidade \cite{hertneck2018}, além de métodos determinísticos baseados em constantes de Lipschitz para assegurar que as discrepâncias entre a inferência da rede e o comando ótimo nominal permaneçam limitadas dentro da margem de robustez do sistema \cite{alsmeier2024}.

% ============================================================
\section{Identificação de Sistemas para Alimentação do NMPC}
\label{sec:fund_sysid}
% ============================================================

A viabilidade prática e o desempenho de qualquer controlador preditivo dependem intrinsecamente da qualidade preditiva do modelo dinâmico utilizado. Em plantas mecatrônicas de aeropêndulo operadas por motores BLDC e hélices, fenômenos como atrito seco nos mancais, histerese aerodinâmica e efeitos de recirculação de ar são de difícil quantificação puramente fenomenológica.

Nesse contexto, as técnicas de Identificação de Sistemas orientadas a dados desempenham papel fundamental. A partir de ensaios dinâmicos em que o sistema é excitado por sinais com riqueza espectral persistente, identificam-se modelos não-lineares — como modelos autorregressivos com entradas exógenas (NARX) ou modelos de caixa cinza (\textit{grey-box}) ajustados via \textit{multiple shooting} —, permitindo alimentar o otimizador preditivo com representações fiéis do envelope operacional real \cite{jiang2022, billings2013}.

% ============================================================
\section{Redes Neurais Multicamadas (MLP)}
\label{sec:fund_mlp}
% ============================================================

Dentre as topologias neurais, as redes \textit{Feedforward} do tipo Perceptron Multicamadas (MLP) são especialmente adequadas para aproximar leis de controle preditivo em razão de seu tempo de execução estritamente determinístico e ausência de laços recorrentes internos. A rede é estruturada em uma camada de entrada que recebe o vetor de estados medidos, históricos passados e a trajetória de referências futuras (\textit{preview}), camadas ocultas densas com ativações ReLU e uma camada de saída linear que fornece o comando escalar de controle ótimo $u^*(k)$.

O treinamento é conduzido sob o paradigma supervisionado com a função de perda do Erro Quadrático Médio (MSE), utilizando algoritmos de otimização estocástica baseados em gradiente, como o Adam. O banco de dados de treinamento é gerado pela resolução offline do problema de controle ótimo em malha fechada cobrindo sistematicamente o espaço de estados de interesse \cite{karg2020, paszke2019}.

% ============================================================
\section{Inferência de Redes Neurais em Hardware Embarcado (TinyML)}
\label{sec:fund_tinyml}
% ============================================================

A implantação de modelos neurais em sistemas mecatrônicos de baixo custo e tempo real restrito insere-se no domínio do \textit{TinyML} (\textit{Tiny Machine Learning}). Em microcontroladores de classe ARM Cortex-M, os recursos de computação são severamente limitados: memórias SRAM da ordem de centenas de kilobytes, memória Flash de alguns megabytes e operação desprovida de sistema operacional de propósito geral.

Para atender a essas restrições, utilizam-se cadeias de ferramentas especializadas. Bibliotecas como o \textit{TensorFlow Lite for Microcontrollers} (TFLM) e extensões proprietárias como o \texttt{X-CUBE-AI} da STMicroelectronics analisam o grafo computacional da rede e o traduzem em código C nativo otimizado, explorando instruções vetoriais da biblioteca CMSIS-NN e unidades de ponto flutuante (FPU). Esse pipeline viabiliza inferências em frações de milissegundo, preservando a determinabilidade temporal requerida para malhas de controle de alta taxa de amostragem \cite{david2021}.

% ============================================================
\section{Trabalhos Correlatos e Lacuna Identificada}
\label{sec:fund_correlatos}
% ============================================================

A convergência entre aprendizado de máquina e controle preditivo tem sido explorada ativamente na literatura recente. \citeonline{jiang2022} aplicaram uma rede neural feedforward operando como modelo dinâmico preditivo acoplado a um otimizador online em um veículo aéreo, enfrentando ainda demandas computacionais consideráveis. Na linha de compensação de perturbações, \citeonline{sanyal2024} propuseram a arquitetura RAMP-Net, que combina Redes Neurais Informadas pela Física (PINNs) para absorver distúrbios de vento em NMPC de drones. Em trabalhos teóricos, \citeonline{li2024} investigaram o Learning-MPC diretamente no grupo de Lie $SO(3)$, exigindo computadores de bordo com GPUs dedicadas para processamento em tempo real.

Nota-se, contudo, que grande parte dos trabalhos foca na aceleração em computadores de bordo pesados ou limita-se à validação puramente simulada. Identifica-se uma expressiva lacuna tecnológica no que concerne ao **pipeline completo de ponta a ponta**: desde a identificação não-linear de uma bancada física real, passando pela síntese offline do NMPC e aproximação neural, até a implantação experimental determinística em microcontrolador de baixo custo (Cortex-M4 a 180\,MHz a 100\,Hz) com diagnóstico e superação do *Sim-to-Real gap*. É a essa lacuna que este trabalho se dedica.
